Quantile trend filtering estimates time-varying conditional quantiles through a check loss and an $\ell_1$ difference penalty, with efficient frequentist estimation available via ADMM \citep{brantley2020baseline}. We develop an adaptive, fully Bayesian extension of this estimator. Under an asymmetric Laplace likelihood and a global Laplace difference prior, the conditional posterior mode at unit likelihood scale and fixed regularization parameter coincides with the frequentist ADMM solution. Assigning priors to both parameters allows the regularization strength and likelihood scale to be inferred jointly with the trend. A three-block Gibbs scheme combines their conjugate updates with a proximal Metropolis-adjusted Langevin (P-MALA) trend update \citep{pereyra2016proximal}, using ADMM to solve the required trend filtering subproblems. The hierarchy uses one global regularization parameter per quantile, without observation-level likelihood augmentation or separate random scales for the trend differences. In the constrained proximal formulation, a fixed banded system permits factorization reuse and linear work in the sample size for a fixed number of inner ADMM iterations.

The framework extends to joint estimation of multiple quantiles under exact non-crossing constraints, enforced through box projections embedded in the ADMM proximal operator, with each quantile updated sequentially by its own P-MALA step. For extreme probability levels, where the near-constant gradient of the asymmetric Laplace likelihood causes gradient-based samplers to drift, we introduce a sentinel quantile strategy: auxiliary quantiles placed beyond the target range whose box constraints confine the target within a dynamically updating interval, preventing drift while preserving exploration.

Simulation studies spanning piecewise constant, varying smoothness, and smooth signals under Gaussian, Beta, and mixed normal noise examine estimation accuracy at interior and extreme quantile levels, with posterior uncertainty summaries obtained from the sampled trends. The results show competitive estimation relative to the Bayesian approach of \citet{onizuka2024fast}. The simpler hierarchy and ADMM computation provide a practical route to extending the frequentist benchmark to joint Bayesian inference.
